\chapter{C的进阶}

C语言之所以能在编程语言界占据这么长时间的统治地位，其直接原因可以被简单的看作是“C是汇编的汇编”。C语言的设计者们在设计C语言时，充分考虑了计算机的底层结构和运行机制，使得C语言能够直接操作内存和硬件资源，从而实现高效的程序执行。

不过，也正因此，C语言的学习者必须深刻的理解计算机的内存模型和运行机制，才能真正的掌握C语言的精髓。C语言的内存模型包括了栈、堆、全局区、代码区等不同的内存区域，每个区域都有其特定的用途和特点。理解这些内存区域的分配和释放机制，对于编写高效、稳定的C程序至关重要。

在C语言的内存模型之上，我们会讲述C语言的最重要特性：指针。它提供了其他语言所没有的、直接操纵内存的能力，使用它可以极大促进编程的效率和灵活性，但同时也带来了更高的复杂性和潜在的风险。指针的使用需要谨慎，理解其工作原理和内存管理是学习C语言的关键，也是后续理解C++语言诸多特性的基础。读者一定要确保理解该内容，而不是背“八股”式的语法，否则后续内容将会变得非常困难。

\section{C的内存模型}

内存可以想象为一条相当长的一维停车场，停车场上有若干块地砖，每一块地砖都标识了一个“位”，可以是0，也可以是1。为了方便起见，通常把8个“地砖”（位）组合在一起，形成一个\textbf{字节}（停车位）。每一个字节都有一个唯一的编号，称为\textbf{地址}。在C语言中，内存的地址是从0开始的整数，通常用十六进制表示： \verb|0x00000000| 、 \verb|0x00000001| 、 \verb|0x00000002| ……。每一个字节都可以存储一个8位的二进制数，也就是一个整数，范围是0到255，即 \mintinline{c}{char} 类型的取值范围。

\subsection{变量的大小和地址}

现在，来了一辆车（一个变量），它需要在停车场上停放（在内存中存储）。为了停放这辆车，我们需要知道它的大小（类型）和它的停车位置（地址）：有的车小，只占一个停车位（\mintinline{c}{char}）；有的车大，占据4个停车位（\mintinline{c}{int}）；有的车更大，占据8个停车位（\mintinline{c}{double}）；还有更大的拖车，车上还载着小车（\mintinline{c}{struct}）。车停好后，我们可以通过车牌号（变量名）来找到它，也可以通过停车位的编号（地址）来找到它。

怎么知道一辆车有多大呢？C语言有算符 \mintinline{c}{sizeof(a)}，它可以告诉我们变量 \mintinline{c}{a} 的大小（单位是字节）。

怎么知道一辆车停在哪呢？C语言有算符 \mintinline{c}{&a}，它可以告诉我们变量 \mintinline{c}{a} 的地址（单位是字节）。

怎么知道某个位置上停了什么车呢？C语言有算符 \mintinline{c}{*p}，它可以告诉我们地址 \mintinline{c}{p} 上停的是什么车（变量的值）。

\subsection{复杂变量的存储}

当然，C并不只有 \verb|char| 、 \verb|int| 、 \verb|double| 这些简单的变量类型。C语言还提供了 \verb|struct| 、 \verb|union| 等复杂类型，它们可以组合多个简单类型，形成更复杂的数据结构。对于这些复杂类型，C语言也有相应的存储规则和对齐要求，以保证数据的正确访问和高效存储。

对于数组，C语言规定：数组的元素必须在内存上连续存储。也就是说，如果我们有一个数组 \verb|int arr[5]| ，那么 \verb|arr[0]| 、 \verb|arr[1]| 、 \verb|arr[2]| 、 \verb|arr[3]| 、 \verb|arr[4]| 这五个元素在内存中是连续排列的，每个元素占据4个字节（假设 \verb|int| 是4字节），它们的地址分别是 \verb|&arr[0]| 、 \verb|&arr[1]| 、 \verb|&arr[2]| 、 \verb|&arr[3]| 、 \verb|&arr[4]| ，并且满足 \verb|&arr[i+1] = &arr[i] + sizeof(int)| 。

对于联合体，C语言规定：联合体的所有成员共享同一块内存区域，它们的地址是相同的。也就是说，如果我们有一个联合体 \verb|union U| ，它有两个成员 \verb|int x| 和 \verb|double y| ，那么 \verb|&u.x| 和 \verb|&u.y| 的地址是相同的，它们都指向联合体的起始位置。联合体的大小是其最大成员的大小。

最难的其实是结构体。结构体的成员在内存中是按照声明的顺序排列的。由于现代CPU以WORD（4字节）为单位访问内存，所以如果同一变量跨越多个字，就会大大降低访问速度。因此，编译器设计者人为规定：\textbf{变量在内存上的起始地址必须是它大小的整数倍}，这样一来，变量就不会跨越多个字了，这在结构体中被称作“内存对齐”。

所以回到停车场比喻上来。假设我们来了一辆拖车，上面载着两个小车，一个是 \verb|int| ，一个是 \verb|char| ，即
\begin{minted}{c}
struct Example {
    char a;      // 1字节
    int b;       // 4字节
};
struct Example example;
int size = sizeof(example); // 8
\end{minted}
停车场管理员（CPU）登记拖车的同时，还要把小车一起登记（改变其在拖车上的位置）；又规定， \verb|int| 必须停在4的倍数号的地块上，而 \verb|char| 可以停在任意地块上。那么，如果这个拖车（结构体）从地块0开始停放，那么 \verb|int| 就必须停在地块4上，而 \verb|char| 只能停在地块0上；中间的地块1、2、3就只能空着，作为填充（padding）字节。这样一来，拖车就占用了地块0到7，共8个地块，而不是5个地块。这个道理和结构体的内存对齐一致，因此上述结构体的大小是8字节。

有了这个理论，我们就可以对结构体的内存布局进行分析了。上述结构体的内存布局如表\cref{tab:struct-offset}所示。

有了这个理论，我们就可以通过调整成员顺序、节省结构体的内存空间。例如以下两个结构体，在语义上是等价的，但在内存布局上却有很大的差异，分别见表\cref{tab:struct-offset2}和表\cref{tab:struct-offset3}。


\begin{minted}{c}
struct Example2 {
    char a;      // 1字节
    double b;    // 8字节
    int c;       // 4字节
};

struct Example3 {
    double b;    // 8字节
    int c;       // 4字节
    char a;      // 1字节
};
\end{minted}

所以说，在定义结构体时，合理安排成员的顺序是非常重要的，学会“见缝插针”，就能最大限度地利用内存空间。一般的，可以把大的成员放在前面，小的成员放在后面，这样可以减少填充字节的数量，从而节省内存空间。

\begin{table}
    \begin{subtable}[t]{0.3\textwidth}
        \centering
        \begin{tabular}{cc}
            \toprule
            偏移量 & 成员 \\
            \midrule
            0 & \verb|char a| \\
            1--3 & 填充字节 \\
            4--7 & \verb|int b| \\
            \bottomrule
        \end{tabular}
        \caption{}
        \label{tab:struct-offset}
    \end{subtable}
    \hfill
    \begin{subtable}[t]{0.3\textwidth}
        \centering
        \begin{tabular}{cc}
            \toprule
            偏移量 & 成员 \\
            \midrule
            0 & \verb|char a| \\
            1--7 & 填充字节 \\
            8--15 & \verb|double b| \\
            16--19 & \verb|int c| \\
            20--23 & 填充字节 \\
            \bottomrule
        \end{tabular}
        \caption{}
        \label{tab:struct-offset2}
    \end{subtable}
    \hfill
    \begin{subtable}[t]{0.3\textwidth}
        \centering
        \begin{tabular}{cc}
            \toprule
            偏移量 & 成员 \\
            \midrule
            0--7 & \verb|double b| \\
            8--11 & \verb|int c| \\
            12 & \verb|char a| \\
            13--15 & 填充字节 \\
            \bottomrule
        \end{tabular}
        \caption{}
        \label{tab:struct-offset3}
    \end{subtable}
    \caption{}
\end{table}

\subsection{CPU怎么执行程序，与函数的内存模型}

为了理解指针的方便性，我们需要了解CPU是如何执行程序的。CPU执行程序的过程可以简单概括为：从内存中读取指令，解码指令，执行指令，写回结果。这个过程是一个循环，CPU不断地从内存中获取下一条指令，并根据指令的内容进行相应的操作。这些指令，可以写成汇编码的形式。

我们不妨写一个非常简单的小程序：
\begin{minted}{c}
int main() {
    int a = 1;
    int b = 2;
    int c = a + b;
    return 0;
}
\end{minted}
这个小程序足够简单，它甚至没有输入输出。

现在，我们来编译一下：
\begin{minted}{bash}
gcc main.c -save-temps -o main
\end{minted}
注意这个 \verb|-save-temps| 参数，它会让编译器在编译时保留中间文件。

于是，比起先前，我们现在的目录下多了几个文件，除了 \verb|main.c| 和 \verb|main| 之外，还有 \verb|main.i| 、 \verb|main.s| 、 \verb|main.o| 。这些文件分别是：
\begin{itemize}
    \item \verb|main.i| ：预处理文件，包含了所有的头文件和宏展开后的代码。
    \item \verb|main.s| ：汇编文件，包含了编译器生成的汇编代码。
    \item \verb|main.o| ：目标文件，包含了编译器生成的机器码。
\end{itemize}
\verb|main.o| 我们是看不懂的， \verb|main.i| 和原代码等价所以也不用看了，我们来看一下 \verb|main.s| ：

\begin{minted}{gas}
	.file	"main.c"
	.text
	.globl	main
	.type	main, @function
main:
.LFB0:
	.cfi_startproc
	pushq	%rbp
	.cfi_def_cfa_offset 16
	.cfi_offset 6, -16
	movq	%rsp, %rbp
	.cfi_def_cfa_register 6
	movl	$1, -12(%rbp)
	movl	$2, -8(%rbp)
	movl	-12(%rbp), %edx
	movl	-8(%rbp), %eax
	addl	%edx, %eax
	movl	%eax, -4(%rbp)
	movl	$0, %eax
	popq	%rbp
	.cfi_def_cfa 7, 8
	ret
	.cfi_endproc
.LFE0:
	.size	main, .-main
	.ident	"GCC: (GNU) 16.1.1 20260728"
	.section	.note.GNU-stack,"",@progbits
\end{minted}
哦天哪，这个也看不懂。不过很快我们就能看懂了。

首先，把没用的东西全都丢掉，比如 \verb|.cfi_*| 、 \verb|.LFB0| 、 \verb|.LFE0| 、 \verb|.size| 、 \verb|.ident| 、 \verb|.section| 这些都是编译器生成的调试信息和元数据，我们不需要关心它们。剩下的就是真正有用的汇编指令了：
\begin{minted}{gas}
main:
    pushq	%rbp
    movq	%rsp, %rbp
    movl	$1, -12(%rbp)
    movl	$2, -8(%rbp)
    movl	-12(%rbp), %edx
    movl	-8(%rbp), %eax
    addl	%edx, %eax
    movl	%eax, -4(%rbp)
    movl	$0, %eax
    popq	%rbp
    ret
\end{minted}
现在的问题是，这里面有一大堆 \verb|%rbp| 这样的东西。它们是什么呢？它们是CPU的寄存器。寄存器是CPU能直接看见、直接操作的数据，它们的访问速度比内存快得多。CPU有很多寄存器，每个寄存器都有自己的名字；一个寄存器是8字节， \verb|%r*| 表示这个数据要占满整个寄存器，而 \verb|%e*| 表示这个数据只占用寄存器的低4字节。比如 \verb|%rbp| 是一个8字节的寄存器，而 \verb|%ebp| 是它的低4字节。

但寄存器相比内存来说，数量极其有限，所以CPU在执行程序时，寄存器只能存放少量的临时数据。为了存放更多的数据，CPU需要使用内存。于是，CPU就有了一个\textbf{栈}（stack）来存放数据。

依然回到停车场比喻上来：栈是一个只有一头能进出的细长停车场，如果后一个车没离开，前一个车就被堵在里面，这就是“后进先出”的特点；一般把最后进去的那辆车叫“栈顶”，最先进去的那辆车叫“栈底”。此外，一般还令栈底的地址比栈顶的地址大（即栈顶是低地址，栈底是高地址），这就意味着栈是向下增长的。CPU在执行程序时，会把一些临时数据（比如函数的参数、局部变量、返回地址等）压入栈中，等到这些数据不再需要时，再从栈中弹出。

但是栈是在内存上的，所以CPU需要一些手段来知道现在车停到哪个位置了。于是，CPU有两个寄存器来管理栈： \verb|%rsp| （stack pointer，栈指针）和 \verb|%rbp| （base pointer，基址指针）。 \verb|%rsp| 总是指向栈顶，也就是最后一辆车停放的位置； \verb|%rbp| 则指向当前函数的栈帧的起始位置，也就是当前函数的局部变量和返回地址所在的位置。特别的，为了防止 \verb|%rsp| 因为少量局部变量而在短时间内多次移动，现代操作系统引入了“红区”机制， \verb|%rsp| 下方的128字节可以被不调用其他函数的函数安全使用，避免了移动；但当局部变量过多、把红区占满之后， \verb|%rsp| 就会被迫移动。

那么，我们就来理解一下这个汇编代码吧。这个函数的栈帧如图\ref{fig:function-stack}所示。
\begin{figure}
    \centering
    \includestandalone[width=0.8\textwidth]{images/function-stack}
    \caption{函数栈帧示意图，这里的每一个块是一个WORD}
    \label{fig:function-stack}
\end{figure}

\begin{itemize}
    \item \mintinline{gas}{pushq %rbp}：把 \verb|%rbp| 的值压入栈中，保存当前函数的栈帧的起始位置。（上一个函数是操作系统用来启动程序的函数，比如说 \verb|_start| ，它的栈帧的起始位置就是 \verb|%rbp| 的值）
    \item \mintinline{gas}{movq %rsp, %rbp}：把 \verb|%rsp| 的值赋给 \verb|%rbp| ，也就是把当前栈顶的位置作为当前函数的栈帧的起始位置。
    \item \mintinline{gas}{movl $1, -12(%rbp)}：把整数1存入 \verb|%rbp-12| 的位置，也就是变量 \verb|a| 的地址。
    \item \mintinline{gas}{movl $2, -8(%rbp)}：把整数2存入 \verb|%rbp-8| 的位置，也就是变量 \verb|b| 的地址。
    \item \mintinline{gas}{movl -12(%rbp), %edx}：把 \verb|%rbp-12| 位置的值（也就是变量 \verb|a| 的值）存入 \verb|%edx| 寄存器。
    \item \mintinline{gas}{movl -8(%rbp), %eax}：把 \verb|%rbp-8| 位置的值（也就是变量 \verb|b| 的值）存入 \verb|%eax| 寄存器。
    \item \mintinline{gas}{addl %edx, %eax}：把 \verb|%edx| 和 \verb|%eax| 的值相加，结果存入 \verb|%eax| 寄存器，也就是计算 \verb|c = a + b| 。
    \item \mintinline{gas}{movl %eax, -4(%rbp)}：把 \verb|%eax| 寄存器的值存入 \verb|%rbp-4| 的位置，也就是变量 \verb|c| 的地址。
    \item \mintinline{gas}{movl $0, %eax}：把整数0存入 \verb|%eax| 寄存器，也就是函数的返回值。
    \item \mintinline{gas}{popq %rbp}：把栈顶的值弹出到 \verb|%rbp| 寄存器，也就是恢复上一个函数的栈帧的起始位置。
    \item \mintinline{gas}{ret}：从当前函数返回，也就是把栈顶的返回地址弹出到 \verb|%rip| 寄存器，继续执行上一个函数的指令。
\end{itemize}

总结一下，汇编里根本就没有什么 a、b、c 这些名字，它只认 \verb|%rbp-12| 、 \verb|%rbp-8| 这样的地址偏移量。换句话说，每一个变量在内存里都有一个独一无二的编号（地址），CPU 就是通过这些编号来读写数据的。

在 C 语言里， \verb|&a| 就能得到 \verb|a| 的地址，它背后的本质就是 RBP 加上一个偏移——这和我们手算 \verb|%rbp-12| 是一模一样的。如果我们能拿到一个变量的地址，并且能像汇编一样通过地址去访问内存……那岂不是可以玩出各种灵活的操作？这恰恰就是 C 语言最强大的武器——指针的起点。

下一节，我们就正式把这个武器握到手里。

\section{指针}\label{sec:pointer}

如果刚才我们讲的内存模型都理解了的话，那么理解指针就不再是一个困难的事情了。

\subsection{什么是指针}

所有教材上都说，“指针是一个存储了地址的变量”，这个定义对于初学者而言比较抽象。但如果继续套用刚才的停车场比喻，我们把指针理解为“停车地块的编号”，就能理解了。

我们知道，对于一个变量，获取其地址的手段是 \mintinline{c}{&a}。这个东西的类型是什么呢？它的类型是“指向a的指针”，如果 \verb|a| 是 \verb|int| 类型，那么 \verb|&a| 的类型就是 \verb|int*| ，也就是“指向整数的指针”。如果 \verb|a| 是 \verb|char| 类型，那么 \verb|&a| 的类型就是 \verb|char*| ，也就是“指向字符的指针”。所以说，指针的类型是由它所指向的变量的类型决定的\footnote{在C语言中，指针实际上并不是一个独立的类型，而是一个复合类型，它由指针本身的类型和它所指向的变量的类型共同决定。但在C++中，我们对指针 \mintinline{cpp}{T*} 进行了严格的定义，使得指针成为了一个独立的类型。我们这里和C++接轨，把指针当作一个独立的类型来看待。}。

如果一个变量占据了多个地址（一辆车占据了多个停车位），那么它的地址就是它所占据的第一个地址（第一个停车位的编号）。所以说，指针指向的是变量的首地址。

对于这种东西，我们通常记作：
\begin{minted}{c}
int a = 1;      
int* p = &a;    // 理解为：p是一个指向整数的指针，它存储了a的地址
int *p = &a;    // 理解为：*p是一个int类型的变量，它存储了a的值
\end{minted}
虽然理解不同，但殊途同归。对于编译器而言，星号写在哪里都无所谓，甚至 \mintinline{c}{int*p = &a;} 和 \mintinline{c}{int * p = &a;} 都是合法的写法。实际写法要和团队的编码规范保持一致。

一个比较坑人的地方是：\mintinline{c}{int* p, q} 中，只有 \verb|p| 是 \verb|int*| 类型，而 \verb|q| 是 \verb|int| 类型。也就是说，星号只和它右边的变量绑定，而不是和类型绑定，所以实际上应该写成 \mintinline{c}{int *p, *q;} 才是两个指针变量。这是C语言的一个历史遗留问题，在C++也没有被解决，容易引起初学者的误解。

为了防止出现错误，C语言建议（注意不是规定）将指针指向的类型和指针本身的类型保持一致。也就是说，如果一个指针是 \verb|int*| 类型，那么它应该指向一个 \verb|int| 类型的变量；如果一个指针是 \verb|char*| 类型，那么它应该指向一个 \verb|char| 类型的变量。这可以避免一些潜在的错误，见下文。

\subsection{初步使用指针}

刚刚我们知道了，定义指针变量相当于拿了一张小纸条，去停车场（内存）把一辆车（变量）所停的停车位的编号（地址）写在上面；使用指针，就是拿着这张小纸条（指针变量）去停车场（内存）按编号（地址）找车（变量）。但在这个过程中，这个编号（内存地址）上停的是什么车（变量），小纸条（指针变量）是不管的。也就是说，指针变量本身并不关心它所指向的变量的内容，它只关心这个变量的地址。

所以说，以下两段代码应该不难理解：
\begin{minted}{c}
    int a = 1;
    int* p = &a; // p存储了a的地址
    printf("%d\n", p); // 0x7ffee3b8c9ac（假设）
    printf("%d\n", a); // 1
    *p = 100; 
    printf("%d\n", p); // 0x7ffee3b8c9ac（和上面一样）
    printf("%d\n", a); // 100
\end{minted}
这个相当于拿着小纸条（指针变量）去停车场（内存）找车（变量），然后把车（变量）的值改成100，但小纸条上还是记的那一串编号。也就是说， \verb|*p = 100| 的本质是把 \verb|a| 的值改成100。

\begin{minted}{c}
int a = 1;
int* p = &a; // p存储了a的地址
printf("%d\n", p); // 0x7ffee3b8c9ac（假设）
printf("%d\n", a); // 1
p = 100;
printf("%d\n", p); // 100
printf("%d\n", a); // 1
printf("%d\n", *p); // Segmentation fault (core dumped)
\end{minted}
这个相当于把小纸条上的编号划去，改成了100。这个过程，对停车场上的车（变量）没有任何影响，所以 \verb|a| 的值还是1。但这次我们试图拿着小纸条找车的时候，发现停车场上根本没有编号为100的车（变量），或者这个车（变量）根本不属于我们，甚至这个编号（地址）根本就不在停车场（内存）上。于是就发生了段错误。

我们发现，刚刚 \verb|p = 100| 之后， \verb|p| 不再指向一个已知的变量，但它确实指向了某段地址（可能极其危险）。这样的指针被称作“悬垂指针”，也俗称“野指针”，对这种指针解引用可能会导致程序崩溃。同理，如果指针初始化时并未立即指向一个已知的变量，也会导致野指针，指向一个垃圾数值。

为了避免这个问题，我们可以让一个指针不指向任何地址，即“空指针”。在C语言中，空指针通常用 \verb|NULL| 来表示。我们可以这样定义一个空指针：
\begin{minted}{c}
int* p = NULL;
\end{minted}
解引用空指针也是不允许的。但至少空指针提供了一个安全的方式来表示“这个指针没有指向任何有效的地址”，而不是随意指向一个可能不存在的地址。

\subsection{利用指针访问数组与指针的算术运算}

刚刚我们知道数组是顺序表结构，即元素在内存中连续存储；换言之，数组的元素在内存中是挨着存放的。所以，我们可以用指针来访问数组中的元素。比如，有一个数组 \verb|int arr[5];| ，它的元素在内存中是连续存储的，我们可以通过指针来访问这些元素：

\begin{minted}{c}
int arr[5] = {1, 2, 3, 4, 5};
int* p = arr; // p指向数组的首元素，也就是arr
printf("%d\n", *p); // 1，访问arr[0]
p++; // p指向下一个元素，也就是arr[1]
printf("%d\n", *p); // 2，访问arr[1]
p += 2; // p指向arr[3]
printf("%d\n", *p); // 4，访问arr[3]
\end{minted}

这里的 \verb|p++| 和 \verb|p += 2| 实际上是指针的算术运算，它们会根据指针所指向的类型来移动地址。比如， \verb|int*| 类型的指针每次加1，实际上是增加了4个字节（假设 \verb|int| 是4字节），因为它要跳过一个 \verb|int| 类型的元素。同理， \verb|double*| 类型的指针每次加1，会增加8个字节（假设 \verb|double| 是8字节）。这也是为什么提倡指针的类型和它所指向的变量的类型保持一致的原因之一：编译器会自动计算偏移量（实际上是乘上一个 \mintinline{c}{sizeof(type)}），从而正确地访问数组中的元素。

当然，通过指针访问数组越界时，也会导致错误，即野指针。

\subsection{利用指针传递地址}

指针可以作为函数的形参，从而实现传递地址的功能。通过传递地址，函数可以直接操作调用者的变量，而不是仅仅操作变量的副本。这种方式被称为“按地址传递”或“传递指针”。

比方说，要交换两个变量。朴素的写法是：
\begin{minted}{c}
void swap(int x, int y){
    int temp = x;
    x = y;
    y = temp;
}

int a = 1, b = 2;
swap(a, b); // 预期 a = 2, b = 1
printf("a = %d, b = %d\n", a, b); // 实际输出 a = 1, b = 2 
\end{minted}
这是因为，在C语言中，函数参数是按值传递的，也就是说，函数接收到的是变量的副本，而不是变量本身。因此，在 \verb|swap| 函数中， \verb|x| 和 \verb|y| 只是 \verb|a| 和 \verb|b| 的副本，对它们的修改不会影响到原来的变量。

但如果写成这样：
\begin{minted}{c}
void swap(int* x, int* y){
    int temp = *x;
    *x = *y;
    *y = temp;
}

// 后面略，总之是符合预期的
\end{minted}
这次， \verb|swap| 函数接收到的是指向 \verb|a| 和 \verb|b| 的指针，也就是它们的地址。通过解引用指针（ \verb|*x| 和 \verb|*y| ），函数可以直接访问和修改 \verb|a| 和 \verb|b| 的值，从而实现交换的功能。这种方式不仅可以实现变量的交换，还可以用于修改数组、结构体等复杂数据类型的内容，极大地提高了函数的灵活性和通用性。

此外，指针还有一个良好的性质：如果传递的实参很大（例如是一个大型结构体或数组），按值传递会导致大量的内存拷贝，影响性能。而通过传递指针，只需要传递地址，避免了不必要的内存开销，提高了程序的效率。一般的，如果不需要修改实参且实参小于8字节时，可以直接按值传递；如果实参大于8字节，或者需要修改实参，则应该按地址传递。

\subsection{指向函数的指针}

指向函数的指针是另一种常见的指针类型。它指向的是函数的地址，而不是变量的地址。通过指向函数的指针，我们可以调用函数，实现函数的动态调用和回调机制。

\begin{minted}{c}
int add(int a, int b) {
    return a + b;
}
int multiply(int a, int b) {
    return a * b;
}

int (*funcPtr)(int, int);
// 将函数指针指向add函数
funcPtr = add;
printf("5 + 3 = %d\n", funcPtr(5, 3)); // 调用add函数
// 将函数指针指向multiply函数
funcPtr = multiply;
printf("5 * 3 = %d\n", funcPtr(5, 3)); // 调用multiply函数
\end{minted}

这样能够让我们在运行时选择要调用的函数，从而实现更灵活的代码结构。

\subsection{指向结构体的指针}

指向结构体的指针和指向基本数据类型的指针类似。例如：
\begin{minted}{c}
struct Point {
    int x;
    int y;
};
struct Point p = {10, 20};
struct Point* ptr = &p; // 定义一个指向结构体的指针
printf("Point p: (%d, %d)\n", ptr->x, ptr->y); // 通过指针访问结构体成员
\end{minted}

\verb|ptr->x| 在这里是一个语法糖，它等价于 \verb|(*ptr).x| 和 \verb|(*ptr).y| ，即先解引用指针，再访问结构体成员。使用指向结构体的指针可以方便地操作结构体数据。

\subsection{常指针和指针常量}

常指针（const pointer）是指针指向的数据不能被修改，但指针本身可以指向其他数据。例如：
\begin{minted}{c}
int a = 1;
int b = 2;
const int* p = &a; // p是一个指向常量的指针，不能通过p修改a的值
// *p = 100; // 错误！
p = &b; // 正确！p可以指向其他数据
\end{minted}    

指针常量（pointer constant）是指针本身不能被修改，但指针指向的数据可以被修改。例如：
\begin{minted}{c}
int a = 1;
int b = 2;
int* const p = &a; // p是一个指针常量，p只能指向a
// p = &b; // 错误！
*p = 100; // 正确！可以通过p修改a的值
\end{minted}

指针常量和常指针可以混合使用，例如：
\begin{minted}{c}
int a = 1;
const int* const p = &a; // p是一个指向常量的指针常量，既不能通过p修改a的值，也不能让p指向其他数据
\end{minted}

\subsection{无类型指针}

有一种特殊的指针：\mintinline{c}{void*}，它被称为无类型指针。

无类型指针可以指向任何类型的数据，但它本身不携带类型信息。由于无类型指针没有类型约束，它不能直接进行解引用操作，也不能进行指针算术运算。在使用此类指针时，必须先将其转换为具体的指针类型，然后才能进行解引用或其他操作。

这个东西有一个妙用：实现最基础的\textbf{运行时多态}。意思是，在运行时，程序可以根据实际情况决定使用哪种数据类型，而不是在编译时就固定下来。

例如
\begin{minted}{c}
void printValue(void* ptr, char type) {
    switch (type) {
        case 'i':
            printf("%d\n", *(int*)ptr);
            break;
        case 'f':
            printf("%f\n", *(float*)ptr);
            break;
        case 'c':
            printf("%c\n", *(char*)ptr);
            break;
        default:
            printf("Unknown type\n");
    }
}
\end{minted}
在这个例子中， \verb|printValue| 函数接受一个无类型指针和一个类型标识符，根据类型标识符的值，函数可以在运行时决定如何解释指针所指向的数据，从而实现对不同类型数据的处理。这种方式在C语言中非常有用，尤其是在实现通用数据结构和库函数时，可以提供更高的灵活性和可扩展性。但这个东西属于较高级的内容，初学者可以先不去深究。

\subsection{基础的直接内存操作}

利用指针，我们可以进行直接的内存操作。这是因为，由编译器自动管理的栈空间是很小的，通常只有几兆字节，而对于一些大型数据结构或者需要动态分配内存的场景，栈空间可能不足以满足需求。为了应对这种情况，C语言提供了动态内存分配的机制，使得程序可以在运行时根据需要申请和释放内存。

\paragraph{堆}

我们先来看看直接内存操作的重要场景：堆。这东西和栈是相对的，特点包括：
\begin{itemize}
    \item 很大：堆的大小通常只受限于系统的可用内存，而不像栈那样有固定的大小限制。
    \item 灵活：堆内存的分配和释放是由程序员手动控制的，而不像栈那样由编译器自动管理。
    \item 缓慢：由于堆内存的分配和释放需要操作系统的参与，因此相对于栈内存来说，堆内存的访问速度较慢。
    \item 任意进出：堆内存的分配和释放可以在程序的任何位置进行，而不像栈那样只有一头能进出。
\end{itemize}

有了堆，我们就可以在程序运行时动态地申请和释放内存，从而实现更灵活的数据结构和算法。C语言提供了几个标准库函数来进行堆内存的管理，包括 \verb|malloc| 、 \verb|calloc| 、 \verb|realloc| 和 \verb|free| 。

\paragraph{基本的动态内存分配}

\begin{itemize}
    \item \mintinline{c}{void* malloc(size_t size)}：分配一块指定大小的内存，返回一个指向这块内存的指针。分配的内存未初始化，内容是未知的。
    \item \mintinline{c}{void* calloc(size_t num, size_t size)}：分配一块指定数量和大小的内存，并将其初始化为零。返回一个指向这块内存的指针。
    \item \mintinline{c}{void* realloc(void* ptr, size_t new_size)}：重新分配一块内存，将原来的内存块调整为新的大小。如果新的大小大于原来的大小，新的内存区域的内容是未定义的。返回一个指向新内存的指针。
    \item \mintinline{c}{free(void* ptr)}：释放之前分配的内存块，指针 \verb|ptr| 必须是通过 \verb|malloc| 、 \verb|calloc| 或 \verb|realloc| 分配的内存块。释放后，指针 \verb|ptr| 仍然指向原来的地址，但该地址的内容不再有效（即野指针）。因此，通常在释放内存后，将指针设置为 \verb|NULL| 是一个良好的编程习惯，以避免悬垂指针的出现。
\end{itemize}

这些东西的使用方法很简单，下面是一个简单的例子：
\begin{minted}{c}
#include <stdio.h>
#include <stdlib.h>

int main() {
    int* arr = (int*)malloc(5 * sizeof(int)); // 分配一个包含5个整数的数组
    if (arr == NULL) { // 检查内存分配是否成功
        fprintf(stderr, "Memory allocation failed\n");
        return 1;
    }

    for (int i = 0; i < 5; i++) {
        arr[i] = i + 1; // 初始化数组
    }

    for (int i = 0; i < 5; i++) {
        printf("%d ", arr[i]); // 输出数组元素
    }
    printf("\n");

    free(arr); // 释放内存
    arr = NULL; // 避免悬垂指针

    return 0;
}
\end{minted}
有的同学可能认为这是脱裤放屁，为什么不直接\mintinline{c}{int arr[5];}就好了呢？这是因为，刚刚这种定义方法定义的 \verb|arr| 是局部变量，在栈上；而如果数组开得非常大，就会导致栈溢出（Stack Overflow），程序崩溃。为了避免这种情况，我们可以把数组放在堆上，这样就不会受到栈大小的限制了。

除此之外，在上一章中我们提到，C标准中的VLA是一个可选特性，所以\mintinline{c}{int arr[n];}这种写法在某些编译器中是不被支持的，例如MSVC。但如果使用malloc来分配内存，就可以在运行时动态地确定数组的大小，从而实现更灵活的内存管理。

堆上还可以放大型结构体，以及链表、树等数据结构。这些数据结构的大小在编译时可能无法确定，或者需要在运行时动态调整大小，因此使用堆内存是非常合适的。

除了这些基本用法之外，还有内存映射等多种高级用法，这些内容属于操作系统的范畴，初学者可以先不去深究。

\paragraph{内存泄漏与双重释放} 

内存泄漏，是指程序在运行过程中，动态分配的内存没有被释放，导致这些内存无法被再次使用，从而浪费了系统资源。内存泄漏可能会导致程序占用越来越多的内存，最终可能导致系统性能下降甚至崩溃。

双重释放是指程序在同一个内存块上多次调用 \verb|free| 函数，这会导致未定义行为，可能引发程序崩溃或安全问题。为了避免这种情况，应该确保每个通过 \verb|malloc| 、 \verb|calloc| 或 \verb|realloc| 分配的内存块只被释放一次，并在释放后将指针设置为 \verb|NULL| ，以防止悬垂指针的出现。

内存泄漏和双重释放都是常见的内存管理问题，程序员在使用动态内存分配时需要格外注意，确保每块内存都被正确地释放，并避免重复释放同一块内存。C++的各种RAII机制都是为了解决这个问题而设计的，但在C语言中，程序员需要手动管理内存，因此需要格外小心。

\section{多文件编程与链接过程}

在实际工程中，如果把所有的代码都写在一个文件中，代码量会非常庞大，难以维护和管理；如果遇到多人协作开发，大家把所有代码都写在一个文件中，冲突和合并时的问题会更加严重。因此，我们需要将代码拆分成多个文件，每个文件负责不同的功能模块，从而提高代码的可维护性和可读性。

C的代码一般被放在头文件 \verb|.h| 和源文件 \verb|.c| 中。在编译过程中，编译器会将所有的源文件分别编译成对应的目标文件 \verb|.o| ，然后链接器会将这些目标文件链接成最终的可执行文件。在这个过程中，编译器和链接器会处理函数和变量的声明与定义，从而确保程序的正确性。变量和函数统称为“符号”，它们在编译和链接过程中扮演着重要的角色。

\subsection{声明和定义}

C语言中，我们声明一个变量 \mintinline{c}{int a;} ，这意味着我们告诉编译器，变量 \verb|a| 是一个整数类型的变量，但并没有为它分配内存空间。这个过程被称为“声明”。而当我们写 \mintinline{c}{int a = 1;} 时，这不仅是声明，还为变量 \verb|a| 分配了内存空间，并初始化为1，这个过程被称为“定义”。

简单变量的声明和定义分得并不开。但一旦变量的类型复杂起来，就会变成：
\begin{minted}{c}
struct Point; // 声明一个结构体类型Point

struct Point { // 定义一个结构体类型Point
    int x;
    int y;
};
\end{minted}
这里的 \verb|struct Point;| 是一个前向声明，它告诉编译器，存在一个名为 \verb|Point| 的结构体类型，但并没有提供它的具体定义。而后面的 \verb|struct Point { ... };| 则是对这个结构体类型的完整定义，包含了它的成员变量。

同理，我们之前说的“函数原型”，也是函数的声明：
\begin{minted}{c}
int add(int, int); // 声明一个函数add，它接受两个整数

int add(int a, int b) { // 定义函数add
    return a + b;
}
\end{minted}

那么声明和定义有什么区别呢？声明只是告诉编译器“有这么一个变量或函数”，而定义则是告诉编译器“这个变量或函数的具体内容是什么”。变量或函数被使用前，一定要先声明它们，而定义可以放在任意位置，只要在链接阶段能够找到它们的定义即可。（变量的定义就是赋值，赋值确实不能放在任意位置，必须在使用时进行赋值。）

在C23中，变量不能被：
\begin{itemize}
    \item 重声明：同一个变量在同一个作用域中不能被多次声明，否则会导致编译错误。
    \item 重定义：同一个变量在同一个作用域中不能被多次定义，否则会导致链接错误。
    \item 冲突声明：同一个变量在同一个作用域中不能被声明为不同的类型，否则会导致编译错误。
\end{itemize}

\subsection{变量的作用域、生命周期}

变量的作用域（scope）是指变量在程序中可见和可访问的范围，生命周期（lifetime）是指变量在内存中存在的时间段。根据变量的定义位置和修饰符，C语言中的变量可以具有不同的作用域和生命周期。
\begin{itemize}
    \item 块作用域（block scope）：在函数内部定义的变量，它们的作用域仅限于所在的代码块（通常是花括号 \verb|{}| 包围的区域）。这些变量在代码块执行时创建，在代码块结束时销毁。
    \item 文件作用域（file scope）：在函数外部定义的变量，它们的作用域是整个源文件，甚至可以被其他文件访问。全局变量和静态变量通常具有文件作用域。
    \item 全局作用域（global scope）：在整个程序中都可见的变量，通常是通过 \verb|extern| 关键字声明的全局变量。它们的生命周期贯穿整个程序运行期间。
\end{itemize}

根据不同的作用域和生命周期，变量可以分为以下几类：
\begin{itemize}
    \item 局部变量（local variable）：在函数内部定义的变量，它们的作用域仅限于函数内部。局部变量在编译时通常不会出现在符号表中，因为它们的生命周期仅限于函数调用期间，编译器会为它们分配栈空间。
    \item 静态变量（static variable）：在函数内部或文件内部定义的变量，但使用了 \verb|static| 关键字修饰。静态变量的作用域是定义它的函数或文件，但它们的生命周期贯穿整个程序运行期间。静态变量在编译时也会出现在符号表中。
    \item 全局变量（global variable）：在函数外部定义的变量，它们的作用域是整个文件，甚至可以被其他文件访问。全局变量在编译时会出现在符号表中，因为它们的生命周期贯穿整个程序运行期间。
\end{itemize}

\subsection{头文件和源文件}

头文件一般放声明，源文件一般放定义。头文件的作用是让编译器知道某个变量或函数的存在，而源文件则提供了它们的具体实现。除此之外，程序必须有且仅有一个 \verb|main| 函数作为程序的入口点。编译器会从 \verb|main| 函数开始执行程序，并根据函数调用关系逐步执行其他函数。

\begin{minted}{c}
// math_utils.h
#ifndef MATH_UTILS_H
#define MATH_UTILS_H
int add(int a, int b);
int subtract(int a, int b);
#endif // MATH_UTILS_H
\end{minted}
\begin{minted}{c}
// math_utils.c
#include "math_utils.h" // 包含头文件
int add(int a, int b) {
    return a + b;
}
int subtract(int a, int b) {
    return a - b;
}
\end{minted}
\begin{minted}{c}
// main.c
#include <stdio.h>
#include "math_utils.h" // 包含头文件
int main() {
    int x = 10;
    int y = 5;
    printf("x + y = %d\n", add(x, y));
    printf("x - y = %d\n", subtract(x, y));
    return 0;
}
\end{minted}

实际编译时，可以使用以下命令将多个源文件编译成一个可执行文件：
\begin{minted}{bash}
gcc main.c math_utils.c -o my_program
\end{minted}
注意头文件不会出现在编译命令中：它们不是“编译单元”（源文件）。

注意 \verb|#include| 指令。该指令的意思是：将后面的文件内容\textbf{直接插入}到当前位置，就像是把文件的内容复制粘贴到这里一样。它后面跟着的如果是尖括号，编译器会优先在标准库中查找；如果是双引号，编译器会优先在当前目录中查找。这类以 \verb|#| 开头的指令被称为“预处理指令”，它们在编译之前由预处理器处理。

类似的，这里的 \verb|#ifndef| 、 \verb|#define| 和 \verb|#endif| 指令是用来防止头文件被重复包含的。它们的意思是：如果没有定义 \verb|MATH_UTILS_H| ，就定义它，并包含头文件的内容；否则，跳过头文件的内容。这种方式被称为“头文件保护”，可以避免重声明、重定义的问题。

\subsection{链接过程}

不同的源文件（编译单元）会生成不同的预处理后代码、汇编码、目标文件，并会被链接到一个可执行文件。前面的预处理、编译、汇编三个过程和单文件编程是一样的，区别在于链接过程。

链接的主要作用是将各个目标文件中的符号（即函数、变量）进行解析、匹配、重定位，并生成最终的可执行文件。

\paragraph{符号表} 每一个目标文件中都会有一个符号表，它记录了该目标文件中声明、定义、引用的符号。会出现在符号表的符号包括：全局变量（包括外部变量）、静态变量、函数。

举一个例子：
\begin{minted}{c}
extern int a;
int b; 
static int c;

int main(){
    int d;
    static int e;
}
\end{minted}
这段代码中，会出现在符号表中的变量有：
\begin{itemize}
    \item \verb|a| ：这是一个外部变量（external variable），它在其他文件中定义，但在当前文件中声明。编译器会在符号表中记录它的名称和类型，以便在链接阶段解析它的定义。
    \item \verb|b| ：这是一个全局变量，它在当前文件中定义。编译器会在符号表中记录它的名称、类型和存储位置，以便在链接阶段正确地分配内存。
    \item \verb|c| ：这是一个静态变量，它在当前文件中定义。编译器会在符号表中记录它的名称、类型和存储位置，但它的作用域仅限于当前文件。
    \item \verb|e| ：这是一个静态局部变量，它在 \verb|main| 函数中定义。编译器会在符号表中记录它的名称、类型和存储位置，但它的作用域仅限于 \verb|main| 函数。
\end{itemize}
而 \verb|d| 这个局部变量不会出现在符号表中，因为它的作用域仅限于 \verb|main| 函数内部，编译器会为它分配栈空间，但不会在符号表中记录它的名称。

函数这东西没有局部的，所以它们都会出现在符号表中。

\paragraph{静态链接}

在链接前，编译器和汇编器会默认所有的声明过的符号都已经被良好的定义了，即使它们在文件外。对于这些外部符号的调用，编译器会生成一个占位符，表示“这里要用某符号，具体位置待定”。此外，由于不知道自己将会在最终程序中的什么位置，编译器会使用相对地址来表示这些符号的引用，即我们看到的所有汇编文件都是从 \verb|0x00000000| 开始的。

随后，链接器会：
\begin{itemize}
    \item 扫描所有符号表，收集所有的符号定义和引用。
    \item 对于每一个符号引用，链接器会查找符号表，找到对应的符号定义，并将占位符替换为实际的地址。
    \item 对于每一个符号定义，链接器会为它分配内存空间，并将其地址记录在符号表中。
    \item 对所有的符号引用进行重定位，即将相对地址转换为绝对地址，确保程序在运行时能够正确地访问这些符号。
    \item 如果找不到某个符号的定义，链接器会报错，提示“未定义的引用”。
\end{itemize}
这种在编译时就把所有东西链接好的做法叫做静态链接。静态链接的优点是程序运行时不需要依赖外部库文件，所有的代码和数据都已经被打包到可执行文件中，运行时加载速度快；缺点是可执行文件体积较大，更新库文件需要重新编译和链接。

\paragraph{动态链接}

动态链接和静态链接的过程大同小异，依然需要扫描符号表、收集符号定义和引用、替换占位符、重定位等步骤。但不同的是：
\begin{description}
    \item[编译时] 编译器在可执行文件中生成一个导入表，记录所需外部函数的信息；
    \item[运行时] 链接器会在程序加载时，根据导入表查找所需的动态链接库（DLL或.so文件），并将这些库加载到内存中；随后，动态链接器会解析所有未静态链接的符号引用，将它们与动态库中的符号定义进行匹配，并将占位符替换为实际的地址；最后，动态链接器会对所有的符号引用进行重定位，确保程序在运行时能够正确地访问这些符号。在这个过程中，动态库的代码一般被编译为位置无关的代码，即便加载到内存中的位置不同，也能正确运行。 
\end{description}
这样做的优势是，不同程序可以共享一份动态库的代码，节省内存空间；同时，更新动态库时，不需要重新编译和链接整个程序，只需要替换动态库文件即可；最后，程序可以在运行时动态加载卸载特定的动态库，实现更灵活的功能扩展和插件机制，如热更新等。缺点是程序启动比静态链接慢，且需要管理依赖，比较不方便。

\subsection{利用标准库}

C标准库头文件一共29个，其中的很多方法可以为我们节省代码量。下面列出一些最常用的标准库头文件及其功能：
\begin{itemize}
    \item \verb|<stdio.h>| ：提供输入输出函数，如 \verb|printf| 、 \verb|scanf| 、 \verb|fopen| 等。
    \item \verb|<stdlib.h>| ：提供通用工具函数，如 \verb|malloc| 、 \verb|free| 、 \verb|exit| 、 \verb|qsort| 等。\verb|qsort| 是一个快速排序函数，它可以对数组进行排序，使用时需要提供一个比较函数。
    \item \verb|<string.h>| ：提供字符串处理函数，如 \verb|strlen| 、 \verb|strcpy| 、 \verb|strcmp| 、 \verb|memcpy| 等，能极大降低字符串操作的复杂度。
    \item \verb|<math.h>| ：提供数学函数，如 \verb|sin| 、 \verb|cos| 、 \verb|sqrt| 、 \verb|pow| 等。
    \item \verb|<time.h>| ：提供时间和日期函数，如 \verb|time| 、 \verb|clock| 、 \verb|difftime| 等。
    \item \verb|<ctype.h>| ：提供字符处理函数，如 \verb|isalpha| 、 \verb|isdigit| 、 \verb|toupper| 、 \verb|tolower| 等。
    \item \verb|stdbool.h| ：提供布尔类型和布尔值，如 \verb|bool| 、 \verb|true| 、 \verb|false| 等。
\end{itemize}